% !TEX root = FlowMatchingOverfitting.tex
%% Pierre Guilbert -- short note skeleton.

\documentclass[11pt,reqno,a4paper]{amsart}

%% ----- Encoding & fonts -----
\usepackage[utf8]{inputenc}
\usepackage[T1]{fontenc}
\usepackage{lmodern}

%% ----- Language -----
\usepackage[english]{babel}

%% ----- Mathematics -----
\usepackage{amsmath,amssymb,amsthm,mathtools}
\usepackage{mathrsfs}

%% ----- Hyperlinks -----
\usepackage{xcolor}
\usepackage[colorlinks=true,linkcolor=blue!60!black,citecolor=blue!60!black,urlcolor=blue!60!black]{hyperref}
\usepackage{orcidlink}

%% ----- Graphics -----
\usepackage{graphicx}
\graphicspath{{../}{./}}
\usepackage{float}    %% provides the [H] placement specifier

%% ----- Lists (compact spacing) -----
\usepackage{enumitem}

\usepackage{placeins}

%% ----- Section heading style (amsart-compatible; titlesec is incompatible
%%       with amsart, so we redefine via the low-level \@startsection macro) -----
\makeatletter
\renewcommand{\section}{\@startsection{section}{1}{0pt}%
  {3.5ex plus 1ex minus .2ex}%
  {1.8ex plus .2ex}%
  {\normalfont\LARGE\bfseries\color{black}}}
\renewcommand{\subsection}{\@startsection{subsection}{2}{0pt}%
  {2.5ex plus .8ex minus .2ex}%
  {1.2ex plus .2ex}%
  {\normalfont\Large\bfseries\color{black}}}
\renewcommand{\subsubsection}{\@startsection{subsubsection}{3}{0pt}%
  {1.8ex plus .6ex minus .2ex}%
  {0.9ex plus .2ex}%
  {\normalfont\large\bfseries\color{black}}}

%% amsart prints author names through \MakeUppercase by default, which renders
%% as "PIERRE GUILBERT". We override \@setauthors to keep the author name in its
%% original (mixed-case) form while preserving amsart's centered layout.
\renewcommand{\@setauthors}{%
  \begingroup
    \def\thanks{\protect\thanks@warning}%
    \trivlist
    \centering\footnotesize \@topsep30\p@\relax
    \advance\@topsep by -\baselineskip
    \item\relax
    \author@andify\authors
    \def\\{\protect\linebreak}%
    \authors
    \par\endtrivlist
  \endgroup
}
\makeatother

%% ----- Theorem environments -----
\theoremstyle{plain}
\newtheorem{theorem}{Theorem}[section]
\newtheorem{proposition}[theorem]{Proposition}
\newtheorem{lemma}[theorem]{Lemma}
\newtheorem{corollary}[theorem]{Corollary}

\theoremstyle{definition}
\newtheorem{definition}[theorem]{Definition}

\theoremstyle{remark}
\newtheorem{remark}[theorem]{Remark}

%% ----- Operators -----
\DeclareMathOperator{\Mat}{Mat}
\DeclareMathOperator{\Hess}{Hess}
\DeclareMathOperator{\diag}{diag}
\DeclareMathOperator{\Ima}{Im}
\DeclareMathOperator{\tr}{tr}
\DeclareMathOperator{\dive}{div}

%% ----- Shortcuts -----
\newcommand{\RR}{\mathbb{R}}
\newcommand{\NN}{\mathbb{N}}
\newcommand{\KL}{D_{\mathrm{KL}}}
\newcommand{\LL}{\mathcal{L}}
\newcommand{\Ppos}{\mathcal{P}_{+}}
\newcommand{\Mman}{\mathcal{M}}
\newcommand{\Tm}{T\mathcal{M}}
\newcommand{\FR}{\mathrm{FR}}
\newcommand{\inner}[1]{\left\langle #1 \right\rangle}

%% Key identities. Use inside a display: \[ \keyformula{...} \]
\newcommand{\keyformula}[1]{\boxed{\,#1\,}}

%% ----- Metadata -----
%% The title concerns the behaviour illustrated here, not every generative model.
\title[Flow Matching Without Double Descent]
{When Generative Models Do Not Exhibit Double Descent: A Flow Matching Example}

\author{Pierre Guilbert\orcidlink{0009-0002-1582-4435}}
\thanks{\textit{Email}: \texttt{spguilbert@gmail.com}}
\email{spguilbert@gmail.com}

\date{\today}

\begin{document}

\begin{abstract}
\end{abstract}

\maketitle

\section{Introduction}
\label{sec:introduction}

Generative modelling is a field of machine learning that aims at
approximating an unknown data distribution, with density
$p_{\mathrm{data}}$, from a finite set of independent samples
$(x_1,\ldots,x_N)\in E^N$, where $E=\RR^d$. Methods include generative
adversarial networks (GANs)~\cite{GoodfellowEtAl2014} and variational
autoencoders~\cite{KingmaWelling2014}.
Diffusion models~\cite{HoEtAl2020,SongEtAl2021} and Flow
Matching~\cite{LipmanEtAl2023} tackle modelling and sampling together by
learning an evolution from a simple initial distribution to the data
distribution. We are interested here in Flow Matching.


\begin{figure}[H]
\centering
\href{run:overfit/figures/spiral_flow_matching_online_975_particles.mp4}{%
\includegraphics[width=\textwidth]{overfit/figures/spiral_online_note_preview.pdf}}
\par\smallskip
{\small\href{run:overfit/figures/spiral_flow_matching_online_975_particles.mp4}{Open the particle video (MP4)}}
\caption{Gaussian-to-spiral transport learnt with fresh training samples.
The video follows 975 particles; click the preview to open the accompanying
MP4. The example is described in Section~\ref{sec:visualisation}.}
\label{fig:online-spiral-video}
\end{figure}

Let $p_{\mathrm{init}}$ be the initial density, and let
$Z\sim p_{\mathrm{init}}$ and $X\sim p_{\mathrm{data}}$ be independent.
We consider the path of random variables
\begin{equation}\label{eq:overfit-interpolation}
X_t=tX+(1-t)Z,\qquad 0\leq t\leq 1.
\end{equation}
For $0<t<1$, its density $p_t$ is the convolution of the data and initial
densities contracted by $t$ and $1-t$, respectively
\cite{Guilbert2026FlowMatching}:
\begin{equation}\label{eq:overfit-convolution}
p_t=
\left[\frac{1}{t^d}p_{\mathrm{data}}\!\left(\frac{\cdot}{t}\right)\right]
*\left[\frac{1}{(1-t)^d}p_{\mathrm{init}}\!\left(\frac{\cdot}{1-t}\right)\right].
\end{equation}
Flow Matching consists in learning a velocity field
$v_\theta:[0,1]\times E\to E$ to transport $p_{\mathrm{init}}$ to
$p_{\mathrm{data}}$. The target velocity $v_t$ and the density path are
linked by the continuity equation~\cite{LipmanEtAl2023,Guilbert2026FlowMatching}:
\begin{equation}\label{eq:overfit-continuity}
\partial_t p_t=-\dive(p_t v_t).
\end{equation}
Sampling from the learnt distribution is then performed using a classical
numerical integration solver, such as Euler or Runge--Kutta, applied to the
learnt velocity vector field with initial samples drawn from
$p_{\mathrm{init}}$.

\section{Double descent and memorisation}
\label{sec:double-descent}

In supervised machine learning, it has been shown that, under some
conditions, increasing model capacity can lead to a \emph{double-descent}
learning curve~\cite{BelkinEtAl2019,NakkiranEtAl2020}. The test error first
decreases, increases near the interpolation threshold, where the model has
enough capacity to fit the training samples, and decreases again for
larger models. Nakkiran et al.~\cite{NakkiranEtAl2020} also observe double
descent with longer training of a fixed ResNet18 on CIFAR-10 with label
noise. Figure~\ref{fig:supervised-double-descent} illustrates both cases.

\begin{figure}[htbp]
\centering
\includegraphics[width=0.8\textwidth,trim=0 0 505.5bp 0,clip]{overfit/figures/references/nakkiran_model_double_descent.png}
\par\smallskip
\includegraphics[width=0.8\textwidth]{overfit/figures/references/nakkiran_epoch_double_descent.png}
\caption{Double descent for ResNet18 models on CIFAR-10 with label noise.
Top: increasing model width. Bottom: increasing training time for fixed
architectures. Panels reproduced from Figures~1 (left) and~9 (left) of
Nakkiran et al.~\cite{NakkiranEtAl2020}.}
\label{fig:supervised-double-descent}
\end{figure}
\FloatBarrier

Recent research on diffusion models reports a different behaviour.
Farghly et al.~\cite{FarghlyEtAl2026} derive, in their analysis, a U-shaped
test loss as a function of model capacity. Urfin et
al.~\cite{UrfinEtAl2026} find that increasing capacity can drive
unregularised models towards memorisation. Along the training-time axis,
Bonnaire et al.~\cite{BonnaireEtAl2025Memorization} identify an initial
generalisation regime followed by a memorisation regime when using finite
training data.

For Flow Matching trained on a fixed dataset, sampling $X$ from the
available points makes the prescribed endpoint the empirical measure
\begin{equation}\label{eq:overfit-empirical-measure}
\widehat P_N=\frac{1}{N}\sum_{i=1}^N\delta_{x_i}.
\end{equation}
Learning this empirical distribution more precisely concentrates generated samples
around individual training points~\cite{LipmanEtAl2023,BambergerEtAl2026}.
Carr\'e du champ Flow Matching proposes a way to regularise this
quality--generalisation tradeoff by introducing anisotropic noise adapted
to the local data geometry~\cite{BambergerEtAl2026}.
We will visualise the transition from generalisation to memorisation on a
small two-dimensional example.

\section{A two-dimensional visualisation}
\label{sec:visualisation}

We propose to visualise this memorisation regime through a two-dimensional
example. We construct $p_{\mathrm{data}}$ by smoothing the spiral
\begin{equation}\label{eq:spiral-curve}
\gamma(s)=hs\bigl(\cos s,\sin s\bigr),\qquad
0\leq s\leq 6\pi,\qquad h=\frac{1.1}{2\pi},
\end{equation}
with a Gaussian kernel of standard deviation $0.16$. The curve parameter
is sampled uniformly. The resulting density is restricted to
$[-3.94,3.94]^2$ and normalised. Figure~\ref{fig:spiral-density} shows this
density and a few samples.

\begin{figure}[htbp]
\centering
\includegraphics[width=0.85\textwidth]{overfit/figures/spiral_density.pdf}
\caption{The spiral density $p_{\mathrm{data}}$ and 150 independent samples.
Colours and contours represent the probability density.}
\label{fig:spiral-density}
\end{figure}

\subsection{Training from an infinite dataset}
\label{subsec:infinite-data}

We first train an MLP with four hidden layers of width 128, SiLU
activations and a sinusoidal time encoding. We use
$p_{\mathrm{init}}=\mathcal N(0,\operatorname{Id})$. At each of the 20,000
training updates, we draw 1,024 fresh samples from $p_{\mathrm{data}}$,
independent Gaussian samples from $p_{\mathrm{init}}$, and times uniformly
in $[0,1]$. Training therefore uses an on-demand stream of samples from the
underlying distribution. The video in Figure~\ref{fig:online-spiral-video}
shows 975 initial Gaussian particles transported by the learnt field,
using 50 fourth-order Runge--Kutta steps.
\FloatBarrier

\subsection{Training from a finite dataset}
\label{subsec:finite-data}

We now draw 1,000 independent samples from the same $p_{\mathrm{data}}$:
500 are used for training and 500 for validation. These samples remain
fixed for the full training. The finite-data model has four hidden layers
of width 256, SiLU activations, the same time encoding, and 64 fixed
spatial Fourier frequencies. Its architecture remains unchanged
throughout this experiment.

We train the network for 100,000 epochs, where one epoch is a full pass
over the 500 training samples, in two batches of 250. Gaussian samples
and uniform times are drawn again at each visit.\footnote{After epoch
30,000, updates use the exact empirical conditional-mean velocity to
reduce gradient noise. This has the same Flow Matching optimum; the
plotted losses keep the original conditional regression definition.}
Figure~\ref{fig:finite-data-losses} shows the training and validation
velocity-regression losses. For evaluation, we use 32 fixed noise/time
draws per data point at every checkpoint.

\begin{figure}[htbp]
\centering
\includegraphics[width=\textwidth]{overfit/fixed_dataset_longer/visualisations_975/figures/fixed_spiral_losses.png}
\caption{Training and validation losses for the fixed 500/500 split.
Left: uniform times in $[0,1]$. Right: times above $0.9$. Vertical marks
identify the three selected checkpoints.}
\label{fig:finite-data-losses}
\end{figure}

The transported initial law at different training regimes is shown in
Figure~\ref{fig:finite-data-transport}: generalisation at epoch 8,500,
the beginning of overfitting at epoch 17,000, and memorisation at epoch
100,000. The three videos follow the same 975 particles sampled from
$p_{\mathrm{init}}$. As training continues, their final positions
concentrate around the fixed training points.

\begin{figure}[htbp]
\centering
\href{run:overfit/fixed_dataset_longer/visualisations_975/figures/fixed_spiral_late_training.mp4}{%
\includegraphics[width=\textwidth]{overfit/fixed_dataset_longer/visualisations_975/figures/fixed_spiral_endpoints.png}}
\par\smallskip
{\small
\href{run:overfit/fixed_dataset_longer/visualisations_975/figures/fixed_spiral_generalisation.mp4}{Generalisation}
\quad
\href{run:overfit/fixed_dataset_longer/visualisations_975/figures/fixed_spiral_memorisation_onset.mp4}{Overfitting begins}
\quad
\href{run:overfit/fixed_dataset_longer/visualisations_975/figures/fixed_spiral_late_training.mp4}{Memorisation}}
\caption{Transported densities at $t=1$ for the three training regimes.
Hollow circles mark the fixed training points. Click the labels to open
the corresponding videos, or the image for the last video.}
\label{fig:finite-data-transport}
\end{figure}
\FloatBarrier

%% ----- References -----
%% Existing bibliography reused verbatim, with the new references appended.
\begin{thebibliography}{99}

\bibitem{NicolaescuCoarea}
L.~I.~Nicolaescu,
\emph{The co-area formula}, lecture notes, revised 2021.
\url{https://www3.nd.edu/~lnicolae/Coarea.pdf}.

\bibitem{LipmanEtAl2023}
Y.~Lipman, R.~T.~Q.~Chen, H.~Ben-Hamu, M.~Nickel, and M.~Le,
\emph{Flow Matching for Generative Modeling}, ICLR, 2023.
\url{https://arxiv.org/abs/2210.02747}.

\bibitem{AlbergoBoffiVandenEijnden2025}
M.~S.~Albergo, N.~M.~Boffi, and E.~Vanden-Eijnden,
\emph{Stochastic Interpolants: A Unifying Framework for Flows and Diffusions},
Journal of Machine Learning Research \textbf{26} (2025), no.~209, 1--80.
\url{https://www.jmlr.org/papers/v26/23-1605.html}.

\bibitem{Otto2001}
F.~Otto,
\textit{The geometry of dissipative evolution equations: the porous medium equation},
Communications in Partial Differential Equations,
\textbf{26} (2001), no.~1--2, 101--174.
\newblock DOI: \href{https://doi.org/10.1081/PDE-100002243}
{10.1081/PDE-100002243}.

\bibitem{BenamouBrenier2000}
J.-D.\ Benamou and Y.\ Brenier,
\textit{A Computational Fluid Mechanics Solution to the Monge--Kantorovich Mass Transfer Problem},
Numerische Mathematik, vol.\ 84, no.\ 3, pp.\ 375--393, 2000.
\newblock DOI: \href{https://doi.org/10.1007/s002110050002}
{10.1007/s002110050002}.

\bibitem{Villani2009}
C.~Villani,
\textit{Optimal Transport: Old and New},
Grundlehren der mathematischen Wissenschaften, vol.~338,
Springer, Berlin, Heidelberg, 2009.
\newblock DOI: \href{https://doi.org/10.1007/978-3-540-71050-9}
{10.1007/978-3-540-71050-9}.

\bibitem{LiuGongLiu2023}
X.~Liu, C.~Gong, and Q.~Liu,
\emph{Flow Straight and Fast: Learning to Generate and Transfer Data with Rectified Flow},
ICLR, 2023. \url{https://arxiv.org/abs/2209.03003}.

\bibitem{EsserEtAl2024}
P.~Esser et al.,
\emph{Scaling Rectified Flow Transformers for High-Resolution Image Synthesis},
ICML, PMLR \textbf{235} (2024), 12606--12633.
\url{https://proceedings.mlr.press/v235/esser24a.html}.

\bibitem{ZhuEtAl2024}
Y.~Zhu, X.~Liu, and Q.~Liu,
\emph{SlimFlow: Training Smaller One-Step Diffusion Models with Rectified Flow},
ECCV, 2024.
\url{https://www.ecva.net/papers/eccv_2024/papers_ECCV/papers/10822.pdf}.

\bibitem{AvrahamiEtAl2025}
O.~Avrahami, O.~Patashnik, O.~Fried, E.~Nemchinov, K.~Aberman,
D.~Lischinski, and D.~Cohen-Or,
\emph{Stable Flow: Vital Layers for Training-Free Image Editing},
CVPR, 2025, 7877--7888.
\url{https://openaccess.thecvf.com/content/CVPR2025/html/Avrahami_Stable_Flow_Vital_Layers_for_Training-Free_Image_Editing_CVPR_2025_paper.html}.

\bibitem{LeEtAl2023}
M.~Le et al.,
\emph{Voicebox: Text-Guided Multilingual Universal Speech Generation at Scale},
NeurIPS \textbf{36} (2023).
\url{https://arxiv.org/abs/2306.15687}.

\bibitem{VyasEtAl2023}
A.~Vyas et al.,
\emph{Audiobox: Unified Audio Generation with Natural Language Prompts},
arXiv:2312.15821, 2023.
\url{https://arxiv.org/abs/2312.15821}.

\bibitem{PolyakEtAl2024}
A.~Polyak et al.,
\emph{Movie Gen: A Cast of Media Foundation Models},
arXiv:2410.13720, 2024.
\url{https://arxiv.org/abs/2410.13720}.

\bibitem{KiEtAl2025}
T.~Ki, D.~Min, and G.~Chae,
\emph{FLOAT: Generative Motion Latent Flow Matching for Audio-driven Talking Portrait},
ICCV, 2025, 14699--14710.
\url{https://openaccess.thecvf.com/content/ICCV2025/html/Ki_FLOAT_Generative_Motion_Latent_Flow_Matching_for_Audio-driven_Talking_Portrait_ICCV_2025_paper.html}.

\bibitem{Bogachev2007}
V.~I.~Bogachev,
\emph{Measure Theory}, Vols.~I--II, Springer, 2007.
\newblock DOI: \href{https://doi.org/10.1007/978-3-540-34514-5}
{10.1007/978-3-540-34514-5}.

\bibitem{AmbrosioGigliSavare2008}
L.\ Ambrosio, N.\ Gigli, and G.\ Savar\'e,
\textit{Gradient Flows: In Metric Spaces and in the Space of Probability Measures},
2nd ed., Lectures in Mathematics ETH Z\"urich,
Birkh\"auser, Basel, 2008.
\newblock DOI: \href{https://doi.org/10.1007/978-3-7643-8722-8}
{10.1007/978-3-7643-8722-8}.

\bibitem{AmbrosioGigli2013}
L.~Ambrosio and N.~Gigli,
\emph{A user's guide to optimal transport},
in \emph{Modelling and Optimisation of Flows on Networks},
Lecture Notes in Mathematics \textbf{2062}, Springer, 2013, 1--155.
\url{https://cvgmt.sns.it/media/doc/paper/195/users_guide-final.pdf}.

\bibitem{Mallat2025}
S.~Mallat,
\emph{Mod\`eles probabilistes par \texorpdfstring{\'energie}{energie} de Gibbs et GAN},
lecture in \emph{G\'en\'eration de donn\'ees en IA par transport et
 d\'ebruitage (2)}, Coll\`ege de France, 2024--2025.
\url{https://www.college-de-france.fr/fr/agenda/cours/generation-de-donnees-en-ia-par-transport-et-debruitage/modeles-probabilistes-par-energie-de-gibbs-et-gan}.

\bibitem{MetropolisEtAl1953}
N.~Metropolis, A.~W.~Rosenbluth, M.~N.~Rosenbluth, A.~H.~Teller,
and E.~Teller,
\emph{Equation of State Calculations by Fast Computing Machines},
The Journal of Chemical Physics \textbf{21} (1953), no.~6, 1087--1092.
\newblock DOI: \href{https://doi.org/10.1063/1.1699114}{10.1063/1.1699114}.

\bibitem{HoEtAl2020}
J.~Ho, A.~Jain, and P.~Abbeel,
\emph{Denoising Diffusion Probabilistic Models},
NeurIPS \textbf{33} (2020), 6840--6851.
\url{https://proceedings.neurips.cc/paper/2020/hash/4c5bcfec8584af0d967f1ab10179ca4b-Abstract.html}.

\bibitem{SongEtAl2021}
Y.~Song, J.~Sohl-Dickstein, D.~P.~Kingma, A.~Kumar, S.~Ermon,
and B.~Poole,
\emph{Score-Based Generative Modeling through Stochastic Differential Equations},
ICLR, 2021.
\url{https://arxiv.org/abs/2011.13456}.

\bibitem{TongEtAl2024}
A.~Tong, K.~Fatras, N.~Malkin, G.~Huguet, Y.~Zhang,
J.~Rector-Brooks, G.~Wolf, and Y.~Bengio,
\emph{Improving and generalizing flow-based generative models with minibatch optimal transport},
Transactions on Machine Learning Research, 2024.
\url{https://arxiv.org/abs/2302.00482}.

\bibitem{Guilbert2026Atlas}
P.~Guilbert,
\emph{Fisher--Rao and Wasserstein--Otto Geometries of Parametric Probability Families: A Comparative Atlas},
preprint, Zenodo, 2026.
\newblock DOI: \href{https://doi.org/10.5281/zenodo.22921933}{10.5281/zenodo.22921933}.



\bibitem{BonnaireEtAl2025Memorization}
T.~Bonnaire, R.~Urfin, G.~Biroli, and M.~M\'ezard,
\emph{Why Diffusion Models Don't Memorize: The Role of Implicit Dynamical Regularization in Training},
Advances in Neural Information Processing Systems \textbf{38} (2025).
\url{https://papers.nips.cc/paper_files/paper/2025/hash/ceb7f3cc876a6dcb15130a645b5a4507-Abstract-Conference.html}.

\bibitem{FarghlyEtAl2026}
T.~Farghly, B.~Dupuis, A.~Durmus, and U.~Simsekli,
\emph{Benign Overfitting Does Not Occur in Diffusion Models},
arXiv:2607.02671, 2026.
\url{https://arxiv.org/abs/2607.02671}.

\bibitem{UrfinEtAl2026}
R.~Urfin, T.~Bonnaire, G.~Biroli, and M.~M\'ezard,
\emph{Double Descent and Malign Overfitting in Diffusion Models},
arXiv:2609.26392, 2026.
\url{https://arxiv.org/abs/2609.26392}.

\bibitem{BambergerEtAl2026}
J.~Bamberger, I.~Jones, D.~Duncan, M.~Bronstein,
P.~Vandergheynst, and A.~Gosztolai,
\emph{Carr\'e du champ flow matching: better quality-generalisation tradeoff in generative models},
ICLR, 2026.
\url{https://proceedings.iclr.cc/paper_files/paper/2026/hash/dbc9f9d486bd1714896111b0fcd78cd1-Abstract-Conference.html}.

\bibitem{Guilbert2026FlowMatching}
P.~Guilbert,
\emph{Flow Matching and Optimal Transport: An Illustrated Introduction to Wasserstein--Otto Geometry},
Zenodo, 2026.
\newblock DOI: \href{https://doi.org/10.5281/zenodo.24084198}{10.5281/zenodo.24084198}.



\bibitem{GoodfellowEtAl2014}
I.~J.~Goodfellow, J.~Pouget-Abadie, M.~Mirza, B.~Xu,
D.~Warde-Farley, S.~Ozair, A.~Courville, and Y.~Bengio,
\emph{Generative Adversarial Nets},
Advances in Neural Information Processing Systems \textbf{27} (2014).
\url{https://proceedings.neurips.cc/paper/2014/hash/f033ed80deb0234979a61f95710dbe25-Abstract.html}.

\bibitem{KingmaWelling2014}
D.~P.~Kingma and M.~Welling,
\emph{Auto-Encoding Variational Bayes}, ICLR, 2014.
\url{https://arxiv.org/abs/1312.6114}.

\bibitem{BelkinEtAl2019}
M.~Belkin, D.~Hsu, S.~Ma, and S.~Mandal,
\emph{Reconciling modern machine learning practice and the bias-variance trade-off},
Proceedings of the National Academy of Sciences, 2019.
\newblock DOI: \href{https://doi.org/10.1073/pnas.1903070116}{10.1073/pnas.1903070116}.

\bibitem{NakkiranEtAl2020}
P.~Nakkiran, G.~Kaplun, Y.~Bansal, T.~Yang, B.~Barak, and I.~Sutskever,
\emph{Deep Double Descent: Where Bigger Models and More Data Hurt},
ICLR, 2020.
\url{https://arxiv.org/abs/1912.02292}.

\end{thebibliography}

\end{document}
